\documentclass[a4paper,22pt]{article}

\usepackage{latexsym}
\usepackage[empty]{fullpage}
\usepackage{titlesec}
\usepackage{marvosym}
\usepackage[usenames,dvipsnames]{color}
\usepackage{verbatim}
\usepackage{enumitem}
\usepackage[pdftex]{hyperref}
\usepackage{fancyhdr}
\usepackage[T2A]{fontenc}
\pagestyle{fancy}
\fancyhf{} % clear all header and footer fields
\fancyfoot{}
\renewcommand{\headrulewidth}{0pt}
\renewcommand{\footrulewidth}{0pt}

\usepackage[english,main=russian]{babel}
\usepackage{ragged2e}
\usepackage{graphicx}

% Adjust margins
\addtolength{\oddsidemargin}{-0.530in}
\addtolength{\evensidemargin}{-0.375in}
\addtolength{\textwidth}{1in}
\addtolength{\topmargin}{-.45in}
\addtolength{\textheight}{1in}

\urlstyle{rm}

\raggedbottom
\raggedright
\setlength{\tabcolsep}{0in}

% Sections formatting
\titleformat{\section}{
  \vspace{-10pt}\scshape\raggedright\large
}{}{0em}{}[\color{black}\titlerule \vspace{-6pt}]

%-------------------------
% Custom commands
\newcommand{\resumeItem}[2]{
  \item\small{
    \textbf{#1}{: #2 \vspace{-2pt}}
  }
}

\newcommand{\resumeItemWithoutTitle}[1]{
  \item\small{
    {\vspace{-2pt}}
  }
}

\newcommand{\resumeSubheading}[4]{
  \vspace{-1pt}\item
    \begin{tabular*}{0.97\textwidth}{l@{\extracolsep{\fill}}r}
      \textbf{#1} & #2 \\
      \textit{#3} & \textit{#4} \\
    \end{tabular*}\vspace{-5pt}
}


\newcommand{\resumeSubItem}[2]{\resumeItem{#1}{#2}\vspace{-3pt}}

\renewcommand{\labelitemii}{$\circ$}

\newcommand{\resumeSubHeadingListStart}{\begin{itemize}[leftmargin=*]}
\newcommand{\resumeSubHeadingListEnd}{\end{itemize}}
\newcommand{\resumeItemListStart}{\begin{itemize}}
\newcommand{\resumeItemListEnd}{\end{itemize}\vspace{-5pt}}

% %-----------------------------
% %%%%%%  CV STARTS HERE  %%%%%%

\begin{document}

%----------HEADING-----------------
\begin{tabular*}{\textwidth}{l@{\extracolsep{\fill}}r}
  \textbf{{\LARGE Дружков Сергей Александрович}} &
  \href{https://github.com/DrSergei}{Github: github.com/DrSergei}\\
  Date of Birth: 23.11.2002\ & City: ~~~~~~~Санкт-Петербург\\
  \href{mailto:serzhdruzhok@gmail.com}{Email: serzhdruzhok@gmail.com}\ & Mobile:~~~~~~+7-900-521-97-42 \\
  \
\end{tabular*}

%-----------EDUCATION-----------------
\section{Education}
  \resumeSubHeadingListStart
    \resumeSubheading
      {СПбГУ (Факультет математики и компьютерных наук)}{Санкт-Петербург, Россия}
      {Прикладная математика и информатика (Современное программирование)}{2021 - 2025}
      {\scriptsize \textit{ \footnotesize{\newline{}\textbf{GPA:} 4.93 (out of 5.00)}}}
      {\scriptsize \textit{ \footnotesize{\newline{}\textbf{Thesis:} Clone detection and automatic refactoring of Lua programs for network devices}}}
    \resumeSubheading
      {ИТМО (Институт прикладных компьютерных наук)}{Санкт-Петербург, Россия}
      {Прикладная математика и информатика (Программирование и искуственный интеллект)\hspace{-50pt}}{2025 - 2027}
      {\scriptsize \textit{ \footnotesize{\newline{}\textbf{GPA:} --- (out of 5.00)}}}
      {\scriptsize \textit{ \footnotesize{\newline{}\textbf{Thesis:} ---}}}
    \resumeSubHeadingListEnd
	    
\vspace{-5pt}
\section{Skills Summary}
    \resumeSubHeadingListStart
        \resumeSubItem{Languages}{~~~~~~~~~~~C/C++, SQL, Go, TypeScript, Lua, Python}
        \resumeSubItem{Technologies}{~~~~~~~~PostgreSQL, SQLite, CMake, Make, Git, Docker, Robot Framework, \LaTeX}
        \resumeSubItem{Communication}{~~~Russian (Native), English (Upper-Intermediate)}
    \resumeSubHeadingListEnd
\vspace{-5pt}

\section{Experience}
  \resumeSubHeadingListStart
    \resumeSubheading{Huawei Russian Research Institute}{}
    {Программист-Исследователь}{Апрель 2022 - настоящее время}
        \resumeItemListStart
        \resumeItem{Исследовательская деятельность}{Исследования в области статического и динамического анализа программ на C/C++, а также языковых инструментов для них.}
        \resumeItem{Рабочая деятельность}{Разработка высокопроизводительных приложений на базе исследований и классических результатов, а также их поддержка и внедрение в существующую экосистему компании.}
      \resumeItemListEnd
\resumeSubHeadingListEnd

%-----------PROJECTS-----------------
\vspace{-5pt}
\section{Projects}
\resumeSubHeadingListStart
\resumeSubItem{Static analysis tool for C/C++}{
\begin{itemize}
\vspace{-5pt}
\item Flow-sensitive interprocedural taint analysis of C/C++ programs.
\vspace{-3pt}
\item Разработка domain-specific проверок и решений для больших проектов.
\vspace{-3pt}
\item Исследование применимости графовых баз данных для задач анализа программ.
\vspace{-3pt}
\end{itemize}
}
\resumeSubItem{Lua minifier}{
\begin{itemize}
\vspace{-5pt}
\item Классические методы уменьшения кода: удаление пробелов, комментариев, пустых строк; переименование переменных и меток; удаление неиспользуемого кода и другие языко-специфичные трансформации кода.
\vspace{-3pt}
\item Возможность поиска и удаления дубликатов фрагментов кода.
\vspace{-3pt}
\item Уменьшение размера кода более чем 2 раза. Время работы < 2 минут на больщой кодовой базе (5 MLOC).
\vspace{-3pt}
\end{itemize}
}
\resumeSubItem{Debugging tools for C/C++}{
\begin{itemize}
\vspace{-5pt}
\item Поддержка отладки C/C++ кода с помощью GDB/LLDB на Linux/Windows/HarmonyOS.
\vspace{-3pt}
\item Полная поддержка Debug Adapter Protocol (DAP) и его внутренних расширений.
\vspace{-3pt}
\item Возможность удаленной отладки оборудования по SSH и Telnet.
\vspace{-3pt}
\item Продвинутые методы отладки (child process debugging, deadlock detector, thread fuzzer, memory debugging, etc.).
\vspace{-3pt}
\item Два MCP-сервера (для VS Code и CLI) на основе DAP для работы с LLM.
\vspace{-3pt}
\item Более 1500 пользователей внутри компании.
\vspace{-3pt}
\end{itemize}
}
\vspace{2pt}
\resumeSubHeadingListEnd

%-----------Awards-----------------
\section{Honors and Awards}
\begin{description}[font=$\bullet$]
\item {Участник заключительного этапа ВсОШ по математике - 2019 (71/130)}
\vspace{-5pt}
\item {Призер заключительного этапа ВсОШ по математике - 2021 (44/166)}
\vspace{-5pt}
\item {Призер заключительного этапа ВсОШ по экономике - 2019 (13/77), 2021 (43/135)}
\vspace{-5pt}
\item {Призер региональной олимпиады по математике для студентов ВУЗов Санкт-Петербурга - 2022 (12/122)}
\vspace{-5pt}
\item {Призер НТО (студенческий трек) по профилю ``Геопространственные цифровые двойники'' - 2023, 2025}
\vspace{-5pt}
\item {Победитель НТО (студенческий трек) по профилю ``Геопространственные цифровые двойники'' - 2024}
\vspace{-5pt}
\item {Призер Я-Профессионал по профилю ``Математика'' - 2025}
\vspace{-5pt}
\item {Победитель конкурса-конференции ``Современные технологии в теории и практике программирования'' - 2025}
\vspace{-5pt}
\item {Победитель сибирской математической олимпиады - 2025 (8/378)}
\vspace{-5pt}
\item {Призер North Countries Universities Mathematical Competition - 2026 (24/149)}
\end{description}


\vspace{-5pt}
\section{Other Experience}
  \resumeSubHeadingListStart
    \resumeSubheading
    {Младший преподаватель в Кировской летней многопредметной школе (ЛМШ)}{Киров, Россия}
    {Подготовка учебных материалов и проведение занятий.}{2021}
    \vspace{-3pt}
    \resumeSubheading
    {Жюри муниципального и областного этапа ВсОШ по математике}{Киров, Россия}
    {Проверка работ участников и выработка критериев.}{2021-2025}
    \vspace{-3pt}
    \resumeSubheading
    {Жюри Санкт-Петербургской математической олимпиады}{Санкт-Петербург, Россия}
    {Проверка решений участников и помощь в организации.}{2022-2026}
    \vspace{-3pt}
    % \resumeSubheading
    % {Участник менторской программы СПбГУ и Nexign}{Санкт-Петербург, Россия}
    % {Разработка карьерного плана для развития в области информационных технологий}{2024}
    % \vspace{-3pt}
\resumeSubHeadingListEnd

\end{document}
